%% sections/02e_hessian_step_function.tex
%% Sezione 2.x — Choice of g(q): Hessian Metric and Physical Motivation
%% Da inserire dopo l'esistente sezione 2.3 (Hessian-Based Step Size Function)
%% come sottosezioni aggiuntive

\subsubsection{Optimality of the Hessian Metric}
\label{sec:optimality}

The spectral radius $\rho(\mathbf{H}(\vq))$ provides a
coordinate-invariant measure of the local curvature of the
potential. We motivate its use as a step function through
three complementary arguments.

\paragraph{1. Orbital mechanics argument.}
The characteristic frequency of the linearised orbit near
$\vq$ is $\omega_{\max} = \sqrt{\rho(\mathbf{H})}$.
The Nyquist condition for resolving this oscillation requires
\begin{equation}
  \Delta t \lesssim \frac{1}{\omega_{\max}}
  = \frac{1}{\sqrt{\rho(\mathbf{H})}},
\end{equation}
which is precisely our step function \eqref{eq:g_function} with
$\varepsilon \sim 1$. This argument shows that $g(\vq) =
\varepsilon/\sqrt{\rho(\mathbf{H})}$ is the \textit{minimum}
step consistent with resolving the local orbital dynamics.

\paragraph{2. Geometric argument.}
The Hessian defines a Riemannian metric on configuration space,
$\langle\delta\vq, \mathbf{H}\delta\vq\rangle$. The eigenvalues
of $\mathbf{H}$ are the principal curvatures of the potential
surface. The largest eigenvalue $\rho(\mathbf{H})$ controls the
fastest growing perturbation — adapting the step to its inverse
ensures that the discrete trajectory deviates from the smooth
Hamiltonian flow by a bounded amount at each step.

\paragraph{3. Comparison with Preto-Tremaine.}
\citet{PretoTremaine1999} derived the optimal step function for
Keplerian motion from a time-transformation analysis:
$g_{\rm PT} \propto r^{3/2}$.
For our choice, $\rho(\mathbf{H}) = 2\mu/r^3$, so
$g = \varepsilon/\sqrt{2\mu}\,r^{3/2} \propto r^{3/2}$:
the Hessian metric recovers the Preto-Tremaine result as a
special case, derived here from a purely geometric argument
without any reference to the Sundman transformation.

\subsubsection{Analytical Hessian for the Full Force Model}
\label{sec:full_hessian}

For the complete force model used in the long-term benchmarks
(Section~\ref{sec:longterm}), the Hessian receives contributions
from each force term.

\paragraph{Keplerian term.}
The Hessian of $V_{\rm Kep} = -\mu/r$ is:
\begin{equation}
  \mathbf{H}_{\rm Kep}^{ij}
  = \frac{\mu}{r^3}\left(\frac{3q_iq_j}{r^2} - \delta_{ij}\right),
  \label{eq:hessian_keplerian}
\end{equation}
with eigenvalues $\{-\mu/r^3,\,-\mu/r^3,\,2\mu/r^3\}$ and
spectral radius $\rho(\mathbf{H}_{\rm Kep}) = 2\mu/r^3$.

\paragraph{$J_2$ term.}
Define
\begin{equation}
  V_{J_2} = -\frac{1}{2}\mu J_2 R_\odot^2
  \frac{3\cos^2\theta - 1}{r^3},
  \qquad \cos\theta = \frac{z}{r}.
\end{equation}
Its Hessian is:
\begin{align}
  A_j &= \frac{30z\,\delta_{jz}}{r^7}
      - \frac{105z^2 q_j}{r^9}
      + \frac{15q_j}{r^7}, \notag \\
  B_j &= \frac{\delta_{jz}}{r^5} - \frac{5z\,q_j}{r^7}, \notag \\
  \mathbf{H}_{J_2}^{ij}
  &= C\left[\left(\frac{15z^2}{r^7} - \frac{3}{r^5}\right)\delta_{ij}
     + q_i A_j - 6\delta_{iz}B_j\right],
  \label{eq:hessian_j2}
\end{align}
where $C = \tfrac{1}{2}\mu J_2 R_\odot^2$ and $z = q_3$.

\paragraph{N-body perturbations.}
For the planetary perturbation from body $k$ at position
$\vq_k$ with gravitational parameter $\mu_k$, define
$\mathbf{u}_k = \vq-\vq_k$ and $r_k = |\mathbf{u}_k|$.
\begin{equation}
  \mathbf{H}_k^{ij}
  = \mu_k\left(
      \frac{3\,u_{k,i}u_{k,j}}{r_k^5}
      - \frac{\delta_{ij}}{r_k^3}
    \right),
  \label{eq:hessian_nbody}
\end{equation}
The full Hessian $\mathbf{H} = \mathbf{H}_{\rm Kep}
+ \mathbf{H}_{J_2} + \sum_k \mathbf{H}_k$ is computed
analytically at each force evaluation, adding negligible
overhead to the existing force model.

The spectral radius is estimated efficiently without full
eigendecomposition via the Gershgorin circle theorem:
\begin{equation}
  \rho(\mathbf{H}) \leq \max_i
  \left|H_{ii}\right| + \sum_{j\neq i} |H_{ij}|,
  \label{eq:gershgorin}
\end{equation}
which provides an upper bound in $\mathcal{O}(9)$ operations.
For the dominant Keplerian term the bound is tight;
for perturbed orbits it overestimates $\rho$ by a factor
of at most 3, giving a conservative (smaller) step size
that ensures safety without compromising accuracy.

\subsubsection{Step Bounds and Practical Implementation}
\label{sec:step_bounds}

Three physical bounds are enforced on the step size:

\paragraph{Local period bound.}
To avoid step-size resonances \citep{HernandezZeebeHadden2022}:
\begin{equation}
  \Delta t \leq \frac{T_{\rm local}}{N_{\rm min}},
  \qquad
  T_{\rm local} = 2\pi\sqrt{\frac{r^3}{\mu}},
  \qquad
  N_{\rm min} = 20.
  \label{eq:period_bound}
\end{equation}

\paragraph{Perihelion bound.}
For $r < 2R_\odot$:
\begin{equation}
  \Delta t \leq \Delta t_{\rm peri} = 1\,\mathrm{min},
  \label{eq:perihelion_bound}
\end{equation}
to handle extreme configurations such as sun-grazing comets.

\paragraph{Yoshida negative sub-step.}
The Yoshida coefficients include $d_2 = w_0 < 0$, which
generates a sub-step backward in pseudo-time. The physical
step for this sub-step is $|\Delta t| = |w_0|\,g\,\Delta s$,
which remains positive and bounded by the same constraints.
The backward pseudo-time step does not imply backward
integration of the equations of motion; it is an artefact of
the composition scheme that cancels in the full DKD sequence
\citep{Yoshida1990}.
